\subsection{比例する高係数は零軸へ縮む}
\label{nhp24:section}

三高同時非負模型の $\delta=sy-tx=0$ について，$x>0$ の全域を排除した．
現在の必要条件は
\[
 x=y=0,\qquad z\ge1,\qquad s<t,\qquad 0\le n\le z.
\]
比例領域全体の不存在が解けたという意味ではない．
公開稿の「比例領域の零軸への帰着」定理と「二整数を固定した零軸の完全有限復元」系を参照する．

再利用する第一の観察は，比例条件では $j=th/s>0$，$F\ge h+1$ なので
$T=h+Fn>0$ から $n\ge0$ となることである．$x\ge1$ なら
\[
 F+1-j=B+(s+t-1)(x-1)+s+sz+y>0
\]
のため $S>0$ から $n=0$ まで従う．
ここでは canonical な向きを保ち，$s,t$ を勝手に交換していない．

$x\ge2$ では，$w=j-h>0$ と置き，差方程式から $b$ を消去して元の pair 残余へ戻す．
$-h^3w^2\mathfrak r$ の $h-1,x-2,w-1,z$ に関する全 $281$ 非零係数が正で，
定数項は $312$ である．$x=1$ は $j=h(v+1)$ と整数性を用いて
$b=h+k$，$k\ge0$（$h\ge2$）へ帰着する．分母を払った元残余は
$k=0$ で正，$k=1$ で負，正次数係数が負なので整数 $k\ge1$ では負である．
$h=1$ の別境界も全係数正の表示で排除する．

残る $x=y=0$ では
\[
 F=b-1+sz,\quad m=F^2+F+1,\quad K=s(F-1)-z+1,\qquad
 t=\frac{mn+K}{b+1}.
\]
この式を元残余に代入して $(b+1)^2$ 倍すると，公開稿の明示四次式
$\mathcal P_{s,z,n}(b)$ を得る．$n\ge z+1$ の領域では
$s=\sigma+1,b=\beta+2,z=\zeta+1,n=\zeta+\nu+2$ とすると，
全 $209$ 非零係数が正で，定数項は $63$ になる．従って $n\le z$．

固定 $s,z$ の全復元は，$n=0,\ldots,z$ の各四次式の全整数根から行う．
$b^4$ の係数は $n^2-zn-1<0$ なので，各式は真の四次式であり，
候補は高々 $4(z+1)$ 個である．整数根定理で全 $b\ge2$ を求めた後，
$t>s$ の整数性，全尺度と六符号，元の全 row/pair と正 weights を検査する．
\nolinkurl{verify_high_d_proportional_v24.py} \texttt{--axis s z} がこの手順を実行する．
これに $b,t$ の探索上限はないが，$s,z$ 自体は未有界である．

次の入口は，四次式を $n=0$，$n=z$，$0<n<z$ に分け，
$b+1\mid mn+K$ と組み合わせた合同障害や整数根の位置を求めることである．
有限個の $s,z$ で模型がないという結果を，全域へ外挿しない．
数学的な新証明の核は全係数証明と $x=1$ の整数の間隙であり，
付属の六組の固定 $(s,z)$ 復元は動作確認である．

また $b+1$ は正整数
\[
 C_n=n\bigl((sz)^2-3sz+3\bigr)+(s^2-1)z-3s+1
\]
を割る．余りが零となる共鳴は $(s,z,n)=(3,1,0),(1,2,2)$ の二つだけで，
前者は $t=s$，後者は元残余 $-b^2+5b+17$ の判別式 $93$ により排除される．
従って $b+1\le C_n\le C_z$ という明示上界も使える．
零余りから自由な $b$ の族を残す必要はない．
